\section{Exception dan try-catch}

\textbf{Exception} adalah mekanisme modern untuk menangani error di PHP menggunakan blok \code{try-catch-finally}. Kode yang mungkin gagal ditempatkan di \code{try}; jika exception dilempar (\textit{thrown}), eksekusi melompat ke \code{catch} yang sesuai. \code{finally} selalu dieksekusi. PHP menyediakan hierarki class exception: \code{Exception} (dasar), \code{RuntimeException}, \code{InvalidArgumentException}, \code{PDOException}, dll. Developer juga dapat membuat custom exception \cite{phpManual}.

\begin{webcode}[language=PHP, caption={Exception handling dan custom exception}]
<?php
class RecordNotFoundException extends RuntimeException {}

function cariProduk(PDO $pdo, int $id): array {
  $stmt = $pdo->prepare("SELECT * FROM produk WHERE id = ?");
  $stmt->execute([$id]);
  $produk = $stmt->fetch();
  if (!$produk) {
    throw new RecordNotFoundException("Produk #{$id} tidak ditemukan.");
  }
  return $produk;
}

try {
  $produk = cariProduk($pdo, 999);
  echo $produk['nama'];
} catch (RecordNotFoundException $e) {
  echo "Not Found: " . $e->getMessage();
} catch (PDOException $e) {
  echo "Database error: " . $e->getMessage();
} finally {
  echo "\nSelesai.\n";
}
?>
\end{webcode}

